\section{The degree-six harmonic realization}
\label{sec:harmonic}

Let \(u_e\in S^2\), \(1\le e\le10\), be the unoriented axes through pairs
of opposite icosahedral faces.  The zonal degree-six harmonics define
\[
 L:\mathbb R^{10}\longrightarrow\mathcal H_6,\qquad
 L(a)(\omega)=\sum_e a_eP_6(u_e\cdot\omega),
\]
where
\[
 P_6(s)=\frac1{16}
 \left(231s^6-315s^4+105s^2-5\right).
\]
As \(A_5\)-modules,
\[
 \mathbb R^{10}\simeq\mathbf1\oplus V_4\oplus V_5,
 \qquad
 \mathcal H_6|_{A_5}\simeq\mathbf1\oplus V_3\oplus V_4\oplus V_5.
\]

The ten axes carry a Petersen adjacency relation.  If \(A\) is its
adjacency matrix, the zonal Gram matrix is expected to be
\[
 K=\frac1{243}(196I+47J-112A).
\]
The Petersen eigenvalues \(3,1,-2\), with multiplicities \(1,5,4\), then
give the positive Gram eigenvalues
\[
 \left(\frac{110}{81}\right)^{(1)},\qquad
 \left(\frac{28}{81}\right)^{(5)},\qquad
 \left(\frac{140}{81}\right)^{(4)}.
\]
Thus \(L\) is injective.  Labeling the axes by pairs
\(\{i,j\}\subset\{1,\ldots,5\}\), the four-space is the Petersen
\((-2)\)-eigenspace
\[
 a_{ij}=y_i+y_j,\qquad \sum_i y_i=0.
\]
These exact assertions form \claimid{PH-1}; C655 must certify them without
floating-point choices and provide an independent replay.

For \(y\in V_4\), put
\[
 F_y(\omega)=\sum_{i<j}(y_i+y_j)P_6(u_{ij}\cdot\omega).
\]
The standard degree-six Gaunt cubic restricts to an \(A_5\)-invariant cubic
on \(V_4\), hence to a scalar multiple of \(\sigma_3\).  The imported exact
calculation predicts
\[
 \frac1{4\pi}\int_{S^2}F_y(\omega)^3\,d\omega
 =
 -\frac{784000}{1247103}\,\sigma_3(y).
\]
The nonzero scalar, its normalization, and the uniqueness input are
\claimid{PH-2}.  They remain gated until C655 supplies the task-owned exact
certificate, independent replay, and primary-source audit.

\paragraph{Order-parameter meaning.}
The Gaunt cubic is the signed third-order invariant used in degree-six
bond-orientational descriptions.  The statement above is narrower than an
identification with the ordinary cubic of a uniform icosahedral neighbour
cage.  The uniform cage lies in the trivial summand; the Clebsch cubic sees
the \(V_4\) component of nonuniform weights on the ten opposite face pairs.
This precise relationship is \claimid{PH-3}.

For a face-axis weight vector \(a\), the projector
\[
 P_{V_4}=\frac{(A-3I)(A-I)}{15}
\]
extracts this component.  A scale-free signed descriptor is then
\[
 C_4(a)=
 \frac{\int_{S^2}L(P_{V_4}a)^3}
 {\left(\int_{S^2}L(P_{V_4}a)^2\right)^{3/2}}.
\]
Possible weights include opposed-face distortion, chemical decoration, or
normal stress.  Whether \(C_4\) predicts rearrangement, yielding, or
nucleation beyond standard quadratic and cubic descriptors is an empirical
question, recorded as \claimid{PH-4}.  No applied performance claim is made
in this paper.
